% §2 Temporal measurement — NEW prose (2026-08-12). Per-axiom motivation in the manner of
% the 2005 IJCV article's §2 (one desideratum per subsection, stated in words first, then
% formalized), assembling the axiom system at the end. The formal statements are shared
% with the blueprint (def:cascade-family, rem:axiom-provenance) and are verbatim
% transcriptions of those nodes, linkage markup stripped.
%
% The subsections motivate; the definition states. A formula appears in a subsection only
% where the subsection has something to derive (the covariance relabelling, the two
% readings of the cascade); everything else is words plus the clause letter.

\section{Temporal measurement}\label{sec:measurement}

A measurement of a temporal signal is an average over an extended region of the immediate
past: a point value is an idealization that no physical aperture realizes, and what an
observer can hold is a \emph{record} of the signal at some resolution, not the signal
itself. A multi-scale front end accordingly maintains a family of such records, indexed by
a scale parameter, and if the front end is \emph{uncommitted} --- built ahead of any
knowledge of what the signals will contain --- the family cannot be derived from a signal
model. It must be derived from general requirements on measurement. This section collects
those requirements. All but one are the requirements of \sscite{fagerstrom2005temporal};
the exception, the cascade property of \S\ref{sec:measurement-cascade}, is the subject of
this paper.

Two asymmetries distinguish time from space. The observer cannot access the future:
measurement must be causal. And the observer cannot return to the past: whatever of the
signal's history is to influence the present output must be carried in the observer's own
state --- the observer must \emph{embody} its past. The two enter the theory differently.
Causality is a condition on the measured values and becomes an axiom directly
(\S\ref{sec:measurement-causality}). Embodiment is a condition on \emph{realizability},
not on the values, and no axiom below mentions it; it returns instead as theorem content,
when \S\ref{sec:signaling} exhibits the state that a covariant observer carries and
\S\ref{sec:jet} reads the temporal derivatives of the signal off that state.

Throughout, signals are integrable transients: $X = \Lone$, with positive cone $X_+$ and
$\int f := \int_{\RR} f(t)\,dt$. The primitive object is the operator carrying the record
at a finer scale to the record at a coarser one,
\[
  \Phi_{x,y} : X \to X, \qquad 0 \le x \le y < \infty,
\]
with the signal itself as the record at scale $0$. That the family is indexed by ordered
\emph{pairs} of scales, rather than by a single scale increment, is deliberate; the reason
is the business of \S\ref{sec:measurement-cascade}. The remaining notation is collected in
\S\ref{sec:preliminaries}.

\subsection{Linearity, positivity, and mass}\label{sec:measurement-linearity}

An uncommitted front end has no preferred signal values: it must treat a superposition of
signals as the superposition of their records, which is linearity, and a measurement must
be stable --- a small signal must leave a small record --- which on $X$ is boundedness of
each $\Phi_{x,y}$ (clause (A1) of Definition~\ref{def:cascade-family} below).

Two further requirements say that smoothing may not create structure of its own. An
intensity signal is nonnegative, and a record that changes sign in response to a
nonnegative signal exhibits structure of the apparatus rather than of the signal;
measurement must preserve the positive cone (A4). And the crudest measurement of a
transient --- its total mass --- is the reading at infinite scale, on which all stages must
already agree: smoothing redistributes mass and neither creates nor absorbs it (A5). A
device that loses mass is admissible physics but different physics --- leakage rather than
measurement --- and is set aside in Remark~\ref{rem:axiom-provenance}.

\subsection{Causality}\label{sec:measurement-causality}

The record at time $t$ may depend on the signal only through its restriction to
$(-\infty, t]$. Formally: if the signal vanishes before some time $t_0$, so does every
record of it (A3). This is the one requirement with no spatial counterpart whatever, and
it is what removes the Gaussian: the Gaussian aperture reads the future, and its causal
truncations destroy the cascade structure discussed below.

\subsection{Covariance}\label{sec:measurement-covariance}

A front end built before its signals has no preferred origin of time and no preferred unit
of time.

No preferred origin: each $\Phi_{x,y}$ commutes with the translations
$(\shift_a f)(t) = f(t-a)$ (A2).

No preferred unit: rescaling the time axis must carry the family into itself. Write
$(\dil_\sigma f)(t) = \sigma^{-1} f(t/\sigma)$ for the mass-preserving dilation. The
subtlety is that nothing so far ties the scale index $x$ to any unit --- the scale axis
carries no intrinsic parametrization --- so the rescaled family can only be required to
reappear \emph{with relabelled indices}: there is an increasing bijection $\Sact_\sigma$
of $[0,\infty)$ with
\[
  \dil_\sigma\,\Phi_{x,y} \;=\; \Phi_{\Sact_\sigma x,\,\Sact_\sigma y}\,\dil_\sigma
  \qquad \text{for all } 0 \le x \le y,\ \sigma > 0,
\]
which is (A8). The relabelling $\Sact_\sigma$ enters as data, not as a choice: the axioms
determine it uniquely, and \S\ref{sec:covariance} shows that a canonical scale
parametrization can be recovered from it. We emphasize that the covariance demanded is
\emph{full}: every $\sigma > 0$, not a discrete ladder $\sigma \in q^{\ZZ}$ of preferred
ratios --- a preferred ratio is a preferred unit in disguise. The weaker discrete form
defines a stratum of its own, taken up in \S\ref{sec:conclusions}.

\subsection{The cascade property}\label{sec:measurement-cascade}

A record at a coarser scale must be computable from the record at any finer scale, without
returning to the raw signal: a measurement of a measurement is a measurement, and the
family closes under composition,
\[
  \Phi_{y,z}\,\Phi_{x,y} \;=\; \Phi_{x,z}, \qquad x \le y \le z,
  \qquad\qquad \Phi_{x,x} = \Id,
\]
which is (A6). This is the axiom the paper changes, and it is worth being precise about
what the classical form asserts. In \sscite{fagerstrom2005temporal}, as throughout the
scale-space literature, the cascade is a one-parameter \emph{semigroup},
\[
  \Phi_\tau\,\Phi_{\tau'} \;=\; \Phi_{\tau + \tau'},
\]
and this single equation contains two separable assumptions. The first is the one just
motivated: composites of measurements are measurements. The second is tacit: that the
composite depends only on the total \emph{amount} of scale crossed --- equivalently, that
the elementary smoothing applied between fine scales equals the one applied between coarse
scales, so that measurement stages are interchangeable. Nothing in the measurement picture
supports the second assumption. A stage that carries a fine record to a slightly coarser
one and a stage that does the same deep in the hierarchy act on different records; there
is no a priori reason the two operations should coincide.

The hemigroup form (A6) keeps the first assumption and drops the second, and nothing else:
scale becomes an ordered interval structure rather than an additive quantity. Two
consequences are immediate. The family must be indexed by pairs, since an increment
$\Phi_{x,y}$ is no longer determined by $y - x$; and the scale axis acquires a gauge
freedom --- any increasing reparametrization of $[0,\infty)$ yields the same family
relabelled --- which is why (A8) had to be stated with the relabelling $\Sact_\sigma$ left
free. The derivation of \S\ref{sec:covariance} closes both books at once: covariance
rigidifies the gauge up to normalization, and the semigroup case is recovered exactly when
stage-interchangeability is reimposed (Corollary~\ref{cor:semigroup-case}).

\subsection{Continuity and nondegeneracy}\label{sec:measurement-continuity}

Measurement degrades gracefully: the record depends continuously on the pair of scales,
jointly, in the norm of $X$ (A7). Together with the diagonal condition
$\Phi_{x,x} = \Id$ this is the two-parameter form of the \emph{extended point} axiom of
\sscite{fagerstrom2005temporal}, $\lim_{\tau \to 0} \phi_\tau = \delta$: as $y \downarrow x$
the record at scale $y$ recovers the record at scale $x$. Finally, the scale axis carries
no idle stretches: strictly coarser means strictly more smoothed, $\Phi_{x,y} \ne \Id$ for
$x < y$ (ND). Nondegeneracy is a normalization rather than a restriction --- a family that
stalls on an interval of scales is the same family on a smaller scale axis.

\subsection{The axiom system}\label{sec:measurement-axioms}

% shared with blueprint def:cascade-family
\begin{definition}[causal cascade measurement family]\label{def:cascade-family}
A family of linear operators
\[
  \Phi_{x,y} : X \to X, \qquad 0 \le x \le y < \infty,
\]
is a \emph{causal cascade measurement family} if it satisfies:
\begin{itemize}
\item \textbf{(A1) Boundedness.} Each $\Phi_{x,y}$ is a bounded operator on $X$.
\item \textbf{(A2) Time-translation covariance.} $\Phi_{x,y}\,\shift_a = \shift_a\,\Phi_{x,y}$
for all $a \in \RR$.
\item \textbf{(A3) Causality.} If $f \in X$ vanishes a.e.\ on $(-\infty, t_0)$, then so does
$\Phi_{x,y} f$.
\item \textbf{(A4) Positivity.} $\Phi_{x,y} X_+ \subseteq X_+$.
\item \textbf{(A5) Unit area.} $\int \Phi_{x,y} f = \int f$ for all $f \in X_+$.
\item \textbf{(A6) Cascade (hemigroup).} $\Phi_{x,x} = \Id$, and
$\Phi_{y,z}\,\Phi_{x,y} = \Phi_{x,z}$ for all $x \le y \le z$.
\item \textbf{(A7) Continuity.} For every $f \in X$, the map $(x,y) \mapsto \Phi_{x,y} f$ is
continuous from $\{0 \le x \le y\}$ to $X$.
\item \textbf{(A8) Scale covariance.} There is a family $(\Sact_\sigma)_{\sigma > 0}$ of
increasing bijections of $[0,\infty)$ such that
\[
  \dil_\sigma\,\Phi_{x,y} = \Phi_{\Sact_\sigma x,\,\Sact_\sigma y}\,\dil_\sigma
  \qquad \text{for all } 0 \le x \le y,\ \sigma > 0.
\]
\item \textbf{(ND) Nondegeneracy.} $\Phi_{x,y} \ne \Id$ whenever $x < y$.
\end{itemize}
\end{definition}

% Verbatim from the blueprint's rem:axiom-provenance (remarks are not statement nodes,
% so this is not a `shared with blueprint` edge; keep the two texts in step by hand).
\begin{remark}[provenance of the axioms]\label{rem:axiom-provenance}
(A1)--(A5) and (A7) are as in \sscite{fagerstrom2005temporal}; the diagonal condition in (A6)
together with (A7) is the two-parameter form of the \emph{extended point} axiom
$\lim_{\tau\to 0}\phi_\tau = \delta$.

The single change is (A6). In \sscite{fagerstrom2005temporal} the cascade reads
$\Phi_\tau \Phi_{\tau'} = \Phi_{\tau + \tau'}$, which additionally assumes that the composite
measurement depends only on the total amount of scale --- that measurement stages are
interchangeable. (A6) drops exactly that assumption, and nothing else: this is the whole
difference between the semigroup axiomatics and the present one, and every result that
follows is a consequence of dropping it.

(A8) is the two-parameter form of scaling covariance; note that an increasing bijection of
$[0,\infty)$ automatically fixes $0$. (ND) excludes idle stretches of the scale axis.

A \emph{passive} variant replaces (A5) by $\int \Phi_{x,y} f \le \int f$ on $X_+$; all results
below hold with sub-probability kernels and an extra killing term, and we do not pursue it.
\end{remark}
